
% =============================================================================
% Project : THE_PHILOSOPHY_OF_VIABILITY
% Book    : Book I — Ontology of Viability
% Book ID : VS.B01
% Part    : Part I — Identity and Manifestation
% Part ID : VS.B01.P01
% Chapter : Chapter 08 — Comparative Philosophical Examination
% Chapter ID : VS.B01.P01.C08
%
% File :
%   4_ASSETS/tables/
%   VS.B01.P01.C08.S13.SS18.P003.TL01.avicenna_claims_and_comparative_limits.tex
%
% Asset :
%   TL01 — Avicenna's Principal Claims and Comparative Limits
%
% Structural Identifier :
%   VS.B01.P01.C08.S13.SS18.P003.TL01
%
% Purpose :
%   To classify the principal Avicennian claims reconstructed
%   in Section 13 and distinguish their logical form from
%   their epistemic or justificatory status in the present
%   comparative examination.
%
% Dependencies :
%   longtable, booktabs, array.
%   These are already loaded by the Book I preamble.
%
% =============================================================================

\begingroup
\small
\renewcommand{\arraystretch}{1.18}
\setlength{\tabcolsep}{4pt}
\setlength{\LTcapwidth}{\linewidth}

% Local table helpers.
\def\colw#1{\dimexpr(\linewidth-8\tabcolsep)*#1/100\relax}

\def\rowsep{\specialrule{0.2pt}{3pt}{3pt}}

\def\tablehead{%
    \toprule
    \textbf{Claim or issue} &
    \textbf{Historical reconstruction} &
    \textbf{Functional or logical form} &
    \textbf{Epistemic or justificatory status} &
    \textbf{Comparative limit for Book I}
    \\
    \midrule
}

\def\grouprow#1{%
    \multicolumn{5}{@{}l}{\footnotesize\textsc{#1}} \\
    \cmidrule{1-5}
    \addlinespace[1pt]
}

\def\EpArg{Argued or Derived Claims}
\def\EpCond{Conditionally Supported Claims}
\def\EpProv{Provisional Constructions}
\def\EpOpen{Unresolved / Open Problems}

\begin{longtable}{@{}
    >{\raggedright\arraybackslash}p{\colw{13}}
    >{\raggedright\arraybackslash}p{\colw{25}}
    >{\raggedright\arraybackslash}p{\colw{14}}
    >{\raggedright\arraybackslash}p{\colw{15}}
    >{\raggedright\arraybackslash}p{\colw{33}}
@{}}

\caption{Avicenna's Principal Claims and Their Comparative Limits}
\label{tab:VS.B01.P01.C08.S13.SS18.P003.TL01}
\\

\tablehead
\endfirsthead

\multicolumn{5}{@{}l}{%
    \footnotesize\textit{
        Table~\thetable\ (continued):
        Avicenna's Principal Claims and Their Comparative Limits
    }
}
\\
\tablehead
\endhead

\multicolumn{5}{@{}r}{%
    \footnotesize\textit{Continued on next page}
}
\\
\endfoot

\bottomrule
\endlastfoot


\grouprow{Metaphysics and the existent}

\textbf{Metaphysics}
&
Metaphysics investigates the existent insofar as it is existent and also establishes the First Principle.
&
Thesis
&
\EpArg
&
The universal subject of metaphysics is not identical with the Necessary Existent or with the totality of causal explanation.
\\
\rowsep

\textbf{The existent}
&
The existent is a primary notion that cannot be defined through a more basic concept.
&
Proposition
&
\EpArg
&
The primacy of the existent does not by itself settle the relation between Being and Identity.
\\
\rowsep

\grouprow{Essence, quiddity, and existence}

\textbf{Essence and existence}
&
The quiddity of a contingent existent does not entail its existence; the relation has modal and causal significance.
&
Thesis
&
\EpArg
&
The distinction is not equivalent to the analytical distinction between Being and Identity.
\\
\rowsep

\textbf{Quiddity in itself}
&
Quiddity considered in itself is neither universal nor particular; these determinations accompany it in modes of existence.
&
Thesis
&
\EpCond
&
The status of quiddity does not establish that it exhausts the ontological constitution of Identity.
\\
\rowsep

\textbf{Thingness and whatness}
&
The reality by virtue of which a thing is what it is is distinguished from the affirmation that the thing exists.
&
Proposition
&
\EpArg
&
The similarity to the formulation of Identity in Book I creates a substantive challenge but not a demonstrated equivalence.
\\
\rowsep

\grouprow{Modality and the First Principle}

\textbf{Necessary Existent}
&
The Necessary Existent is uncaused and necessary in itself.
&
Thesis
&
\EpArg
&
The First Principle is not equivalent to Being as such, Identity, or First Determination.
\\
\rowsep

\textbf{Uniqueness of the Necessary Existent}
&
The Necessary Existent is unique; two independently necessary existents of the same absolute status are not possible within the argument.
&
Proposition
&
\EpArg
&
The uniqueness of the First Principle does not establish that every existent has exactly one identity.
\\
\rowsep

\textbf{Possible existent}
&
A contingent existent is possible in itself and does not contain a sufficient reason for its own existence.
&
Proposition
&
\EpArg
&
Possibility in itself is not a statement of epistemic uncertainty.
\\
\rowsep

\textbf{Necessity through another}
&
A possible existent becomes necessary through its complete cause while remaining possible in itself.
&
Proposition
&
\EpArg
&
Causal necessitation is not equivalent to persistence, invariance, or identity preservation.
\\
\rowsep

\grouprow{Unity, causality, and cosmology}

\textbf{Unity and plurality}
&
Unity is a primary metaphysical notion, but its senses and applications must be distinguished.
&
Proposition
&
\EpArg
&
Unity is not identical with Identity, Singularity, or the unity of an aggregate.
\\
\rowsep

\textbf{Causal series}
&
Caused existents and intermediate causes depend upon a First Principle in the order of causal explanation.
&
Argument
&
\EpCond
&
The argument must not be reduced to a claim about the impossibility of every form of temporal infinity.
\\
\rowsep

\textbf{Emanation}
&
Dependent existents proceed from the First through an ordered causal structure.
&
Thesis
&
\EpCond
&
Emanation is not equivalent to Expression or Manifestation.
\\
\rowsep

\textbf{Cosmological order}
&
Separate intelligences, celestial souls and bodies, and sublunary existents are related through a mediated order.
&
Proposition
&
\EpCond
&
The cosmological order is not a map of the proposed realms and hubs of Book I.
\\
\rowsep

\grouprow{Individuation and the comparative problem}

\textbf{Individuation}
&
Avicenna's discussions involve several possible individuative factors, whose roles differ across contexts and kinds of entities.
&
Open Problem
&
\EpOpen
&
No single principle of individuation is adopted here as a universal criterion of Identity.
\\
\rowsep

\textbf{Quiddity and Identity}
&
The relation between the reality of a thing and the ontological status of Identity requires further analysis.
&
Open Problem
&
\EpOpen
&
Quiddity must not be renamed Identity; its possible constitutive role remains open.
\\
\rowsep

\textbf{Causal dependence and Identity}
&
The causal explanation of the existence of a dependent existent may be relevant to its constitution, but the relation requires independent argument.
&
Hypothesis
&
\EpProv
&
Causal dependence does not automatically provide a complete account of the Identity of the effect.
\\
\rowsep

\textbf{Book I response}
&
Book I distinguishes Being, Identity, quiddity, individuation, and causal dependence without treating them as interchangeable.
&
Thesis
&
\EpProv
&
The positive explanatory role of Identity must be developed beyond a merely negative non-reduction principle.
\\
\rowsep

\textbf{Relation between the systems}
&
The two frameworks address overlapping problems while preserving distinct conceptual and explanatory architectures.
&
Thesis
&
\EpCond
&
The relation is critical proximity without conceptual equivalence or direct doctrinal continuity.
\\

\end{longtable}

\endgroup

% =============================================================================
% END OF TABLE TL01
% =============================================================================
